\section{Evaluación constitutiva y realización dimensional}
\label{iii:sec:evaluacion}

La evaluación reúne ahora las construcciones anteriores en su orden efectivo.
Las expansiones regionales de $\pi$, $\varphi$ y $e$ proceden del refinamiento
nonádico; el registro $K$ conserva sus doce posiciones. Su lectura periódica
determina la sección posicional de $\alpha$. Sobre ella actúan el carácter
de acción y el retorno regional, después el par angular y, finalmente, la
respuesta constitutiva. Esta evaluación es posterior a la selección de las
reglas: ninguna de sus cifras se utiliza para elegir un canal.

\subsection{Evaluación de la cadena interna}

Para reproducir la tabla se emplean las expansiones regionales de mil
decimales generadas por el certificado nonádico y el entero $N_K$ de
\eqref{eq:k-numerador}. La lectura utilizada es
\[
 \alpha=\pi+e-\varphi-4-\frac{N_K}{1000^{12}-1}.
\]
La sección finita de doce bloques y la sección prolongada tienen sus
definiciones y su diferencia exacta en el capítulo del registro. Aquí se
mantiene fija la segunda durante toda la evaluación. No se intercambian
lecturas al aumentar la precisión, ni se sustituye el operador analítico
de la segunda vía por una coincidencia decimal con la primera.

\begin{longtable}{@{}ll@{}}
\toprule
Objeto & Evaluación decimal (24 cifras significativas)\\
\midrule\endhead
$\alpha$ & $0.00729735256928380099728511\ldots$ \\
$H_5$ & $1.05457181764615446962582\ldots$ \\
$\eta_{\rm ret}$ & $1.05457181691503906113369\ldots$ \\
$A_{\deg}$ & $7.29735256928380099728511\ldots$ \\
$C^*_{\deg}$ & $2.12373833896864572426195\ldots$ \\
$q_+$ & $0.848377942576404374945737\ldots$ \\
$q_-$ & $0.913660151150557285295519\ldots$ \\
$r_+$ & $0.999999628548249363178695\ldots$ \\
$r_-$ & $0.999724137189518364131517\ldots$ \\
$\widehat\varepsilon$ & $0.999999257096636702760442\ldots$ \\
$\widehat\mu$ & $0.999448350479326935089982\ldots$ \\
$\widehat Z$ & $0.999724508538937215455554\ldots$ \\
$\widehat c$ & $1.00027631048615752610639\ldots$ \\
$\gamma$ & $0.274007672694459274747255\ldots$ \\
$G_{\rm Cat}$ & $0.915965594177219015054604\ldots$ \\
$\rho$ & $0.267949192431122706472554\ldots$ \\
\bottomrule
\end{longtable}


El cálculo se realiza con 120 cifras de trabajo; se muestran 24 cifras
significativas para facilitar la comparación entre magnitudes de distintos
órdenes. Las 120 cifras son una precisión de evaluación, no una incertidumbre
experimental ni, por sí solas, un certificado de error mediante intervalos.
El paquete conserva las entradas generadas, sus huellas, las operaciones y
los residuos de las identidades comprobadas. Las pruebas algebraicas del
cuerpo establecen esas identidades con independencia de esta tabla.

El significado de las cuatro últimas coordenadas constitutivas se conserva
al evaluar: $\widehat\varepsilon$ y $\widehat\mu$ son las respuestas
cuadráticas por hoja, $\widehat Z$ es su orientación relativa y
$\widehat c$ es el inverso de su producto. La cercanía de algunas cifras a
la unidad expresa el pequeño defecto del cociente $3$--$4$ sobre estos
canales. Su origen se ve exactamente en
\[
 1-f(s)=\frac{s^3}{1+s+s^2+s^3}.
\]
La fórmula permite estudiar la respuesta antes de atribuirle una magnitud
dimensional: el signo del defecto, su orden cúbico y su monotonía son
consecuencias del mismo operador.

\subsection{Coordenatizaciones de unidades y covariancia}

Sean $\mathcal L_S$, $\mathcal L_Q$ y $\mathcal L_C$ las rectas de
acción, carga y velocidad. Una coordenatización física positiva evalúa las secciones
$u_S$, $u_Q$ y $u_C$ en unidades de acción, carga y velocidad. La realización
\eqref{iii:eq:dimensions} produce entonces una impedancia, una velocidad,
una permeabilidad y una permitividad, con sus dimensiones respectivas.
Las coordenadas de la tabla anterior permanecen adimensionales.

\begin{proposition}[Covariancia de la realización constitutiva]
Si las secciones de referencia cambian mediante
$u_S'=a u_S$, $u_Q'=b u_Q$, $u_C'=d u_C$, con $a,b,d>0$, las
coordenadas que representan las mismas magnitudes satisfacen
\[
 \widehat Z'=a^{-1}b^2\widehat Z,\qquad
 \widehat c'=d^{-1}\widehat c,\qquad
 \widehat\mu'=a^{-1}db^2\widehat\mu,\qquad
 \widehat\varepsilon'=adb^{-2}\widehat\varepsilon.
\]
Las relaciones $\widehat\mu'\widehat\varepsilon'=(\widehat c')^{-2}$
y $\widehat\mu'/\widehat\varepsilon'=(\widehat Z')^2$ se conservan.
\end{proposition}
\begin{proof}
Se iguala cada magnitud de \eqref{iii:eq:dimensions} a su expresión
en la nueva base y se despeja su coordenada. Los factores del producto
son $d^2$, y los del cociente son $a^{-2}b^4$, justamente los requeridos
por las coordenadas transformadas de $c$ y $Z$.
\end{proof}

Esta covariancia distingue dos actos: construir una sección y expresar
esa sección en unidades elegidas. El teorema espectral fija las coordenadas
en la normalización interna declarada. Una comparación SI debe exhibir
también la evaluación de las tres secciones de referencia. Ajustarlas
mediante las magnitudes que se desea reproducir sería una calibración de
la coordenatización, no una verificación independiente de aquellas magnitudes.

Para reconstruir los canales desde coordenadas expresadas en la coordenatización
primada se deshace ese cambio antes de utilizar la raíz cúbica:
\begin{equation}
 \widehat Z=ab^{-2}\widehat Z',\quad
 \widehat c=d\widehat c',\quad
 \widehat\mu=ad^{-1}b^{-2}\widehat\mu',\quad
 \widehat\varepsilon=a^{-1}d^{-1}b^2\widehat\varepsilon'.
 \label{iii:eq:undo-units}
\end{equation}
Por ejemplo, las entradas de la inversa espectral son
$r_+=(\widehat Z\widehat c)^{-1/2}$ y
$r_-=(\widehat Z/\widehat c)^{1/2}$, calculadas después de
\eqref{iii:eq:undo-units}. Introducir directamente las coordenadas primadas
en esas expresiones cambiaría la normalización del operador y podría
sacarlas del intervalo $(3/4,1)$. La covariancia dimensional conserva
las magnitudes físicas; la inversión cúbica reconstruye el transporte
en la coordenatización interna en que fue definido.

\subsection{Controles de identidad y alcance de la comparación}

El contraste reproducible se organiza por relaciones. En primer lugar,
$\widehat\mu\widehat\varepsilon\widehat c^2=1$ y
$\widehat\mu/(\widehat\varepsilon\widehat Z^2)=1$ comprueban la
consistencia constitutiva. En segundo lugar, la inversión cúbica recupera
los argumentos y las coordenadas angulares de partida. Finalmente, la
evaluación de Barbero por el funcional angular y por el cálculo funcional
de la respuesta coincide por la identidad probada en el capítulo anterior.
El intercambio de hojas debe invertir la contribución orientada y mantener
el modo común; este control detecta una pérdida de etiquetado que los
productos escalares no revelarían.

La comparación de la constante de Planck reducida con las tablas de
referencia se conserva en el desarrollo antecedente de acción. Para las
cuatro magnitudes del vacío, la presente tabla es una evaluación de sus
coordenadas internas. No se presenta como una tabla de coincidencias SI:
esa afirmación adicional requiere una coordenatización numérica independiente para
$u_S,u_Q,u_C$. El resultado constitutivo demostrado, incluida su inversión,
permanece completo en su dominio y conserva explícitamente la interfaz
dimensional necesaria para realizar esa comparación.
